This thesis investigated whether neuroevolution can produce cost-to-go heuristic functions
for Weighted A* search on Boxoban that are substantially more parameter-efficient than a
conventional gradient-trained CNN baseline. Three approaches were implemented and compared
on 294 held-out levels: a fixed convolutional network with 2,031,809 parameters, a
NEAT-evolved feed-forward network, and a CoDeepNEAT hybrid that combines evolutionary
architecture search with gradient-descent weight training.

The CNN baseline established a strong reference point, achieving a 97.3\% solve rate.
NEAT produced a heuristic of just 24 parameters --- less than 0.002\% of the CNN's
parameter count --- and achieved an 84.0\% solve rate. CoDeepNEAT, operating with
32,000 to 170,000 parameters, achieved an 88.8\% solve rate. Both neuroevolved
heuristics are meaningfully weaker than the CNN baseline, but the scale of the parameter
reduction relative to the modest drop in solve rate suggests that neuroevolution is a
viable strategy for heuristic compression.

\subsection{Future Work}

Several directions merit further investigation. Longer NEAT runs with larger populations
may find topologies with greater expressive capacity without abandoning the parameter
efficiency that makes the approach attractive. CoDeepNEAT could be extended to search
over convolutional module types rather than dense-dropout blocks, which may better suit
the spatial structure of the Sokoban grid. More broadly, the question of whether a
neuroevolved heuristic can fully close the gap with the CNN baseline --- or whether the
CNN's inductive structure provides an irreplaceable advantage --- remains open.